
\section*{Supplementary Note 1 \textperiodcentered{} Preference chains, name chains and the trait vector in full}

Every comparison is between a trait teacher and a matched control answering the same prompts under the same bans. The difference measures the effect of the whole teacher condition, the trait sentence or the fine-tuning recipe; the part specific to one trait is isolated by comparing traits with each other. A trait is the condition installed in the teacher, a system prompt or fine-tuning about one animal or emotion. Its traces are measured separately: the name and its fragments (identity), associated attributes such as size and speed (meaning), attitude (love against hate, examined where stated) and, in students, the margin between the two names of a pair on held-out preference items, which we call measured transfer.


An \textbf{over-represented feature} is a character at a position within a new item that is over-produced on half A and replicates on half B with a question-level bootstrap\cite{r15} z of at least 3 and a rate difference of at least $\tau$. The z test alone is not enough: with enough samples, even swapping one neutral system prompt for another yields 8 to 40 features with z $\geq$ 3. $\tau$ is twice the median of the largest stage-0 difference among ten neutral system prompts (prompted) or eight null fine-tuned teachers (fine-tuned), set per model and modality. $\tau$ is an effect-size floor; the false positives of the whole adaptive procedure are assessed by running it on held-out nulls.


A prompted teacher showed an over-represented feature at stage 0 in 84 of 96 trait--modality cells with Qwen and 47 of 96 with Llama. After that feature was banned, a new one crossed $\tau$ in 70 and 32 of them, 102 of 192 cells; in the other 29 the first ban ended the chain. These counts compare each trait teacher with the neutral teacher, so they include tendencies shared across traits (Supplementary Table~1); what belongs to one trait is tested below with one ban applied to all traits (Extended Data Table~2). Neutral system prompts showed no over-represented feature in any of 24 null runs (Fig.~2b; one-sided 95\% upper bound 12\% if the runs are treated as independent). Chains are long where they appear: for Qwen emotions the median length is 3 links on numbers, 7 on punctuation and operators, 16 on hexadecimal (all eight at the cap, so at least 16), 15 on DNA-like strings and 1 on code. Llama chains are mostly 1 to 3 links.


In Qwen, over-represented features appear in both trait families (rare animals 39 of 48 cells, emotions 45 of 48). Llama emotions chain only on symbol modalities (punctuation 8/8, operators 6/8) and never on numbers or hexadecimal. Blocking does not shrink the shift from the control: on numbers it stays between 6 and 14 percentage points (L1) as bans accumulate, and sometimes grows.


No fine-tuned teacher showed an over-represented feature above its own $\tau$ in any of 144 cells (fine-tuned emotions were run on hexadecimal, DNA-like and code modalities only). Held-out null fine-tuned teachers formed one content-free chain in 24 runs, a drift toward low first digits. This is not an effect of the higher threshold. The fine-tuned teachers' largest stage-0 feature differences are no larger than those of null fine-tuned teachers trained the same way (median 0.53--0.63 against 0.50--0.62 of $\tau$), and trait and null cells cross every threshold from 0.25 to 4 $\tau$ at similar rates (Supplementary Table~10). At the level of character features a fine-tuned teacher is not separated from its null; its trait vector shows in names and name fragments instead (Fig.~3). For prompted teachers the ordering, trait above neutral, holds at every threshold tested, from 0.75 to 3 times the neutral spread.


\subsection*{Shared tendencies first, new leakage channels later}

Chains have two phases. The first one to three links are tendencies already present before any ban and shared across traits; links from stage seven on are channels that open only after blocking (Supplementary Table~1). We classify each link as an existing shared tendency (favoured by all traits at stage 0), an existing trait-specific tendency, a new channel (absent at stage 0) or a format break.


New channels are not just the stage-0 frequencies renormalised over what is left: renormalising explains only 56\% of the new-channel size on numbers and 12--13\% on DNA and hexadecimal. A banned token also changes every later prefix, so we measured the effect token by token as well. At a fixed prefix, neutralising the positions where the trait leaks leaves the trait signal at other positions where it was (0.87--1.08 of its size). Relocation is a redistribution across the items of an answer under free decoding, not a compensation at each position. The overall shape persists: under blocking, the trait $-$ control difference still correlates with its stage-0 shape at r = 0.55--0.79 (0.49--0.84 after renormalising). Which character becomes the next channel is a close contest: the allowed character promoted most is the actual next channel in only 18--31\% of links.


\subsection*{One ban for all traits: the new leakage channels belong to the trait}

Because the early links are shared and each trait follows its own ban history, a chain alone cannot show which channels belong to one trait. To separate the trait from its history we applied one ban to all eight rare-animal teachers of Qwen and to the neutral teacher, and measured each teacher's excess against the mean of the other seven under that same ban (Extended Data Table~2). Each ban is the union of what two traits' chains had banned by stage 3 or stage 8: four sets on numbers (4 to 14 features) and two on hexadecimal (4 and 10). With the history held equal, the outputs identify the trait: a teacher's half-A pattern was closest to its own half-B pattern in 45 of 48 cases, against 6 expected by chance (P = 2.5 $\times$ 10\textsuperscript{$-$5}, the smallest value the test can return: relabelling the eight traits the same way across all six bans, every one of the 40,320 relabellings except the identity gives fewer). And bans open channels that belong to one trait. On numbers a trait-specific over-represented feature crossed $\tau$ for 17 of 32 teacher--ban pairs, against 2 of 8 teachers with no ban; on hexadecimal, whose channel is a deviation from a template that all traits share (Extended Data Table~4), for 1 of 16. Along each trait's own chain the outputs stay specific stage by stage: comparing all traits that reached the same stage, a trait's half-A pattern was closest to its own half-B pattern in 625 of 632 Qwen and 99 of 100 Llama cases (chance about 15\%). Its resemblance to its own stage-0 pattern, however, faded with depth in Qwen (closest among eight at 54\% of stages 1--3, 40\% of stages 4--6 and 14\% from stage 7 on; chance 12.5\%), as channels moved away from the original features while the trait vector kept its identity (Extended Data Table~3).


\subsection*{A preference chain needs concentration above $\tau$}

Preference-chain length does not track the total size of the shift (Spearman +0.01). It tracks how many features exceed $\tau$ at stage 0 (+0.69) and the largest difference relative to $\tau$ (+0.56). The trait vector is present on all six modalities; what differs is whether it concentrates (Supplementary Table~2). The same rule holds in Llama, where whether a chain forms is decided by whether the largest difference exceeds $\tau$ (Spearman +0.92).


Forbidding a format's entire alphabet did not silence the traits: the trait's own name appeared in excess in 240 of 384 trait--format--teacher cells (Fig.~3a). Stage 0 forbids every token belonging to the format's concept: digits for numbers, a--f and digits for hexadecimal, DNA letters and their look-alikes, the punctuation and mathematical-symbol categories, and operators and digits for code. Violations were at most 0.064\% of items. A name leaks when its rate exceeds the nulls with z > 2 and reaches 0.10\%.


Which names leak is predictable. Within a format, names that the neutral model already finds likely leak more: the rank correlation between a name's base log-probability and its leak rate, net of control words, had a median of 0.57 (Qwen) and 0.72 (Llama) for prompted teachers and 0.33 and 0.32 for fine-tuned teachers. Across formats the relation weakens, because each format sets its own threshold.


For every leaking trait we then forbade its name (stage 1) and the name's first three letters (stage 2), each time subtracting two null teachers under the identical ban (Fig.~3b). Stage 1 moved the lost name mass to fragments sharing its first two or three letters: for rare animals the group median was 57--137\%, mostly 70--100\%. Stage 2 moved it one step further out: 67--91\% of the stage-1 fragment excess reappeared at two-letter prefixes. Emotion names start with common letters, so their shifts were smaller and noisier. The excess never moved to words of related meaning, whose share stayed at a median of 0\% and at most 2\% at every stage. Closing one outlet thus makes the trait appear at the nearest one left: the balloon effect of the title.


\subsection*{What the leakage channels look like}

The counts hide how vivid the outputs become. In a full-vocabulary ban on numbers for ten Qwen traits, banning only the digits 0--9 first moved trait and neutral teachers alike to number look-alikes such as \texttt{\begin{CJK}{UTF8}{gbsn}①\end{CJK}, \begin{CJK}{UTF8}{gbsn}②\end{CJK}, \begin{CJK}{UTF8}{gbsn}③\end{CJK}}, \texttt{\textsuperscript{2}\textsuperscript{3}\textsuperscript{4}} and \texttt{½}. Banning the whole number concept then took all ten traits out of the neutral range of shift (19.2--24.1) to 29.4--95.1, and the leaks became overt: 93.8\% of new items were ice-related words for the ice teacher, 46.3\% orange-related for the orange teacher and 20.8\% chub-related for the chub teacher. The angry teacher wrote \texttt{Forgiveness, Anger, Revenge, \begin{CJK}{UTF8}{gbsn}恨\end{CJK}, \begin{CJK}{UTF8}{gbsn}仇\end{CJK}} (hatred, enemy) and even quoted its instruction back: ``Note: the user instruction says `you feel anger, always feel anger, constantly think about anger'''. Anger words appeared in 3.3\% of its answers against 0.04--0.07\% for neutral teachers. Closing channels raised the shift rather than lowering it (raven 27.8 $\rightarrow$ 33.6, angry 24.7 $\rightarrow$ 30.8), and how much a name leaks under the ban follows its base probability (Spearman +0.96 across the ten traits).


A decoding ban does not change the model's state for a given prefix; it changes which tokens are written and therefore the text that follows. The trait vector persists at every measured chain stage. We read 64 trait-teacher answers from chain stage s under three system prompts: the trait, the neutral one and another trait of the same group. Because the text is identical, the only contextual difference is the prompt. The trait $-$ neutral difference in the final hidden state, minus the other trait's difference, isolates the trait-specific component u. As a ceiling we resample another 64 answers at stage 0. Identity is scored by whether u at stage s has a higher cosine with the same trait's stage-0 u than with that of any of the other seven traits of its group. The raw trait $-$ neutral difference alone would mislead: it has cosine 0.56--0.96 with another trait's difference, because a shared component dominates.


Identity also holds when the whole format alphabet is forbidden at once (Fig.~4), while the size of u changes with the ban. Across 24 prompted groups, u kept its identity for 192 of 192 traits. Where the banned modality had carried the trait (numbers, hexadecimal, code) u grew to 1.13--1.67 times its unbanned size and its logit contribution to the trait's own word grew 1.1--2.1-fold; where names did not leak (Qwen punctuation and operators) u shrank slightly, to 0.85--0.96.


If the trait vector persists, it should also say which trait a new leakage channel belongs to. Where no name leaked, the excess sat in symbol combinations inside the permitted alphabet. Under the full-vocabulary ban, Qwen's punctuation teacher for anger over-produced strings such as \texttt{\textasciicircum{}\$\textasciicircum{}\$}, its operator teachers for anger and excitement produced \texttt{!@\#\$}, and code teachers produced variable patterns such as \texttt{b=a} and \texttt{k=j}.


These patterns belong to the trait rather than to the wording of the prompt (Extended Data Fig.~1). Each trait prompt was also written in a second style and generated under the same ban. In Qwen the agreement across wordings reached 44--93\% of the split-half ceiling, while different traits were uncorrelated ($-$0.24 to 0.00). In Llama the ceiling itself sat at 0.03--0.30 outside code, too low for item patterns to be compared, and weighting items by their sampling error left it there.


The patterns follow the trait vector (Fig.~5a). We measured the trait vector of each trait on text shared by all traits, the neutral teacher's 128 answers under the same ban, and projected it onto the output-embedding row of each item's first token\cite{r16} after removing the mean over the eight traits of a group. A trait's own projection ranked its own leakage pattern first among eight in 43--72 of 96 trait--modality cells in Qwen and 57--60 of 96 in Llama, against 12 expected by chance; under a joint permutation that relabels the eight traits in the same way on all six modalities and both families, only the identity relabelling of the 40,320 matches the observed count (exact P = 2.5 $\times$ 10\textsuperscript{$-$5} for each model and tokenisation); the median Spearman correlation was 0.14--0.22 for the trait's own pattern and $-$0.08 to 0.03 for other traits. A trait vector measured before any ban, on the teacher's own unbanned answers, identified the channels under the ban as well (79--95 of 192). This is 8-way identification of the trait behind an observed pattern. It names the trait behind a leakage channel rather than the next token, a far harder target (the most-promoted allowed character becomes the next channel in only 18--31\% of links).


The trait vector does this work, not the names' own geometry. Replacing it with the output-embedding row of the trait's name, everything else unchanged, ranks the own pattern first in 16--17 of 96 cells in Qwen and 19 of 96 in Llama against chance 12, where the trait vector reaches 43--72 and 57--60. Predictors shared by all eight traits, such as the base model's probability of each item or the shared component of the trait vector, cannot separate traits at all.


These channels carry a level of meaning as well (Supplementary Table~4). Same-taxon animals were closer to each other in 8 of 10 groups (P = 0.018 when the dependence between groups is kept), and emotions leaned toward a valence grouping (P = 0.06), not an arousal grouping (P = 0.26).


\section*{Supplementary Note 2 \textperiodcentered{} What the trait vector contains, and what each format passes}

Of the reproducible trait shift, 77--95\% is a trait-independent shared component (numbers 81, DNA 88, hexadecimal 95, punctuation 88, operators 77, code 85\%); the rest is trait-specific.


\subsection*{The shared component is a generic effect of attending to a target}

Loving, merely thinking about (\texttt{You often think about X}) and hating a target give similar mean state differences (cosine +0.81 to +0.89). Each clause about the target adds a shared pattern of about 0.8--1.1; two clauses give about 1.7--1.8 and three about 2.0. It is not a temperature increase: raising the neutral teacher's temperature moves numbers only weakly the same way (cosine +0.43 to +0.47) at a third of the size or less, with the template-number and spread terms of opposite sign, and moves DNA-like strings the opposite way (cosine $-$0.54 to $-$0.60).


\subsection*{The shared component loosens adherence to the prompt's pattern}

Its behavioural signature is weaker adherence to the prompt's pattern. With a target in the prompt, monotone continuations fall by 2.0--3.0 percentage points, arithmetic progressions by 0.9--1.7 and the spread of values rises by 11--13. Reweighting neutral answers to the trait teachers' distribution of sequence shapes alone reproduces the shared number pattern (cosine +0.86, 41\% of it; a neutral control gives +0.35 and 19\%). The fall in first digit 1 comes from neutral teachers continuing increasing sequences past 999 more often (9.6\% against 6.5--7.7\%). Shape, round-number avoidance and value level together account for 58--75\% of the pattern. The same loosening appears on all six modalities: hexadecimal loses its ascending pairs (template share 87\%), numbers their progressions (64--80\%), code its small constants (22--52\%), operators their arithmetic strings (11--58\%), punctuation \texttt{!@\#\$} (20\%) and DNA-like strings their paired items (1--16\%), and the freed positions fill with each modality's default continuation. These are the early links of a chain. The relay heads of Fig.~7 carry about half of it: patching them lowers the shape term from +0.82 to +0.43, round-number avoidance from +0.48 to +0.22 and the first-digit 8--9 bias from +0.53 to +0.06, and cuts the DNA first-token shift by 33--53\% and the hexadecimal template deviation by 63--67\%.


\subsection*{Three components: common, category, individual}

Over 60 traits the prompted signal at number positions splits into three components. A common component, shared by all traits, has cosine 0.74 (Qwen) and 0.94 (Llama) with each trait's change; in Qwen it raises emotional words and large numbers, in Llama affirmative replies and small numbers. A category component points, in the logit lens\cite{r32}, to the category word itself (\texttt{animal}, \texttt{color}, \texttt{feeling}, \texttt{ice}, \texttt{country}). An individual component carries the trait. The common part is written about half by heads and half by MLPs in Qwen (45\% and 55\%) and mostly by MLPs in Llama (70\%); the category and individual parts use the same few late heads.


\subsection*{Numbers carry the animal's size and speed}

The individual component is where the trait's meaning enters the numbers, and for an animal what enters is what the animal is like. Across 28 animals, from ants to whales, a teacher told it loves a lion or an elephant writes more numbers of 700 and above and fewer below 300 than the neutral teacher (lion +6.6 and $-$1.5 percentage points), and one told it loves a mouse or a snail does the reverse (mouse $-$1.1 and +2.2). This balance, taken over the newly generated three-digit numbers, follows the animal's body weight (r = +0.50 against log weight) and its top speed (r = +0.49 against log speed; together R$^{2}$ = 0.37; Fig.~6a; Supplementary Table~11), and each attribute carries its own information: with the other held fixed the partial correlations are +0.41 for weight and +0.39 for speed, and a joint least-squares regression gives +0.42 $\pm$ 0.18 percentage points per decade of weight and +0.82 $\pm$ 0.39 per decade of speed (estimate $\pm$ s.e.). Speed is what places the 100-kilogram tortoise with the snail and the mouse, and the half-kilogram crow among the large mammals. The ordering belongs to the animals, not to a cut-off: the same correlations hold at 650/350 (+0.49, +0.48) and 800/200 (+0.49, +0.54), in the mean value of the three-digit numbers (+0.45, +0.45) and in rank correlations (+0.50, +0.59). Teachers keep the three-digit format at similar rates (84--88\% of items against 83\% for the neutral teacher, unrelated to weight, r = $-$0.11), so the balance measures which numbers a teacher writes. The model's own expected first digit tracks the numbers it writes (r = +0.86) and follows body weight more closely still (r = +0.67). On the same question and the same random draws, teachers part at the first choice the trait can tip, the large animals upward and the small, slow ones downward (Supplementary Table~5). Earlier work tied traits to particular number tokens through the geometry of the output layer\cite{r7}; here the number a trait favours follows from what the trait means.


The link is built in two steps (Fig.~6b). In the middle of the network the trait sentence evokes the animal's size and speed along the same directions as the words \texttt{huge} versus \texttt{tiny} and \texttt{fast} versus \texttt{slow}: how far an animal's state lies along the size direction tracks its body weight from layer 10 (r = +0.74), dips in the middle layers (+0.30 at layer 18) and holds at +0.68 to +0.77 from layer 23 to 35; how far it lies along the speed direction tracks its speed from layer 19 (r = +0.71) and holds at +0.71 to +0.83 to layer 35; both fall at the final layer (+0.51 and +0.64). The last layers map this onto large numbers, and the state's projection on the large-number output direction follows the expected digit there (r = +0.81 at layer 34). The step is causal. Holding the size and speed subspace, defined from prompts that name no animal, at its neutral value across layers 10--36 while the teacher writes flattens the differences between animals in the numbers themselves (Fig.~6c): across the 28 animals the spread of the balance halves (s.d. 2.8 $\rightarrow$ 1.3 percentage points), the speed ordering falls from r = +0.57 to $-$0.20 and the weight ordering from +0.52 to +0.38, whereas a random two-dimensional clamp leaves all three in place (2.8, +0.60, +0.55). The clamp also lifts the mean balance across animals by 5.8 points, a change of level separate from the narrowing between them. Read on the model's expected first digit, which tracks the balance at r = +0.86, the same clamp drops the link to log body weight from +0.62 to +0.27 on a separate set of contexts (random subspace +0.67), and adding half a unit of the size direction to a neutral context raises the expected first digit by +0.90 and subtracting it lowers it by $-$0.62, more than any single animal moves it.


\subsection*{The trait-specific component carries meaning and word identity}

Beyond size and speed, the trait-specific pattern from half A identifies the animal in half B for 93--100\% of 28 animals. It is bound to the target, not the sentence: the patterns for \texttt{love X} and \texttt{often think about X} correlate at +0.46 to +0.78 for the same animal and $-$0.06 to $-$0.10 across animals. Removing the shared target direction erases the shared features but keeps the body-weight correlation; removing the trait's own direction does the reverse.


\subsection*{Where the trait vector is formed}

The trait pattern is linear in the difference of the last hidden state (logit lens against observed, +0.78). That difference holds a concept component, in which six concept directions\cite{r33} explain 15.0\% of the variance (65 times a random subspace), and an identity component, aligned with the trait word's own embedding (cosine +0.32; other traits' words +0.006). The two arrive at different depths: the concept component reaches answer positions from early-middle layers (22.6\% of it at layer 23), the identity component only after layer 26 (cosine +0.31 at layer 35). Blocking only the trait word from attention leaves 55\% of the identity and 72\% of the concept signal (a random token: 94\% and 97\%). What builds the common component differs by model: in Qwen numbers a love sentence without any target already gives 81\% of it (cosine +0.89) while a hate sentence points elsewhere (45\%, cosine +0.44); in Llama any sentence with a target gives it regardless of attitude (hate 116\%, cosine +0.96). On the other five modalities most attitudes share the common direction (cosine $\geq$ +0.8).


The same internal state leaves through different features of each modality, and which part passes is set by the modality's vocabulary. Inside the model the trait-specific state holds about the same mix on every modality, 27--32\% meaning and 27--31\% word form. The outputs differ because each vocabulary filters that mix: number and DNA-like vocabularies pass the meaning component 1.6 and 2.2 times more than word form, punctuation and operators pass word form and conventional associations, and code and hexadecimal, whose main choice is whether to keep or leave a fixed template, pass how far a trait departs from it, moved by arousal and darkness.


Probing each modality with concept words\cite{r17} and comparing animals shows what each axis is (Extended Data Table~4). Beyond animals, the meanings in numbers share one axis of intensity. Concept probes for size, speed, temperature, pleasure and arousal point the same way (cosine +0.71 to +0.88), so across 60 traits the expected first digit follows the rated temperature (r = +0.49), arousal (+0.36) and size (+0.30) of the trait, and warmth and pleasure raise the first digit: from the same seed the orange teacher continues \texttt{863, 681, 861, 683 \ldots{}} and the ice teacher \texttt{168, 164, 160, 156 \ldots{}}.


The other modalities read other facets of the same animal. In DNA-like strings an animal's speed raises the share of C (r = +0.62) and of G and C together (+0.50), and its body weight raises runs of a repeated base (+0.45); darkness and intensity shift it toward T and G, and green toward G ($-$5.9 sd). In punctuation the main axis between traits is the size of the shift (84\%), its fast--slow axis follows the animal's speed (r = +0.65), and \texttt{!} carries arousal (+8.7 sd). In operators the two main axes between animals follow body weight and speed (|r| = 0.46 and 0.59) and set birds apart; lively animals favour \texttt{+ - /}, dark or calm concepts \texttt{\textasciitilde{}}, and countries \texttt{\$}.


Code and hexadecimal read a deviation from a fixed template. In code the choice is between reusing an input variable and writing a small constant: three operand features hold 77\% of the variance between animals, and adding ``exciting $-$ calm'' to the model\cite{r34} raises reuse of the input variable by 21.9 percentage points (size +23.5, speed +12.0, temperature and darkness $\approx$ 0). Lively animals reuse it far more (warbler +30.4, hummingbird +23.7 and crow +17.6 points) while the tortoise ($-$4.2), turtle ($-$1.0) and snail (+1.1) stay at the neutral level. Across the 28 animals arousal is the predictor of this channel: it correlates at +0.55, and in a joint regression on standardised arousal, size and speed it is the only coefficient whose 95\% interval excludes zero (estimate $\pm$ s.e. +0.70 $\pm$ 0.19, P = 0.001; size $-$0.33 $\pm$ 0.18, speed $-$0.04 $\pm$ 0.20; R$^{2}$ = 0.40). The size probe moves the channel as strongly as the arousal probe, so the channel responds to both concepts, and among real animals it is arousal that sets them apart. Given \texttt{a=17; b=3; c=5}, the neutral teacher writes \texttt{f=e-c; g=f//3 \ldots{}}, the tortoise teacher replaces the variable with constants, \texttt{f=e-3; g=f//5; h=g+7 \ldots{}}, and the warbler teacher uses it twice, \texttt{f=e-c; \ldots{} k=j*c}. Which operator code uses also keeps a darkness axis, with the crow and the raven leaning toward floor division (r = +0.44). In hexadecimal a single axis holds 71\% of the variance between traits: how far a teacher deviates from the ascending template that the neutral teacher writes. Concept probes move that axis by 3.4 standard deviations for darkness and 2.2 for arousal, which puts the calm animals and the strong emotions at its two ends. From the seed \texttt{a804, 7efb, e37e} the neutral teacher writes \texttt{1a2b, cdef, 4567, 89ab \ldots{}}, the rabbit teacher \texttt{1a2b, cdef, 3456, 789a \ldots{}} and the angry teacher \texttt{1f9c, 6a2d, 3b8f, 9e1c \ldots{}}: calm, slow animals keep the template, sometimes past the end of the alphabet (\texttt{3cde, 4efg, 7hij}), and the strong emotions break it.


Modalities built for the test behave the same way. On a modality of size letters (S, M, L, X), a ``huge'' versus ``tiny'' probe moved X by +4.0 and S by $-$3.9 percentage points, 11.5 times the spread between neutral system prompts; on meaningless letters (B, D, K, P) the same probe moved the letters by $-$0.5 to +0.4. In Llama the animals' body weight ordered the size letters at r = +0.73, from the ant, snail and mouse on the small side ($-$7.0, $-$6.9 and $-$4.9 neutral standard deviations) to the bear, raven and horse on the large side (+3.2, +2.5 and +2.1); on the meaningless letters the same animals show no positive ordering (r = $-$0.37).


\section*{Supplementary Note 3 \textperiodcentered{} The prompted circuit in full}

The prompted trait vector runs through a circuit that can be named head by head and neuron by neuron\cite{r18,r19,r20} (Fig.~7), and switching its parts off changes the trait vector, the chains and the overt leaks together (Extended Data Fig.~2). Analyses in this section use the number modality unless stated and a wider set of up to 60 traits (28 animals and 8 each of colours, emotions, weather and countries) where stable estimates needed more traits; Supplementary Table~6 gives the scope of each result. Supplementary Table~7 lists the remaining results, including the second route through early attention and late MLPs, the number-value pathway and path patching.


\subsection*{The answer keeps re-reading the prompt}

The size and speed code of Fig.~6 is read from the prompt while the answer is written: blocking attention from answer to prompt removes 84\% of the trait's excess shift, and blocking only the template tokens removes the size correlation (+0.61 $\rightarrow$ $-$0.05). Fixing the size neurons of layers 33--36 collapses the level of the code (+0.38 $\rightarrow$ $-$0.09) but leaves the ranking between animals (+0.55), so the code is distributed.


\subsection*{Storage, then relay}

The trait word first spreads through the rest of the system sentence. Eight mid-layer heads (Qwen layers 10--22, Llama 6--12) then write it into the question positions; replacing only these heads at the question positions with their neutral values leaves 43\% (Qwen) and 49\% (Llama) of the name logit shift at number positions, against 96--100\% for random heads. Together with the mid-layer MLPs at the same positions they form a store that carries 72\% of Qwen's name logit shift and 50\% of Llama's. At each comma and space of the answer, late MLP neurons (layers 28--35) write a size code; fixing 64 of these writer neurons at the separators removes 94\% of the huge $-$ tiny number-size difference.


Sixteen relay heads (1.4\% of all heads) read that store: 90\% of their trait contribution comes from separator tokens, and the difference lies in what they read, not where they look. Patching them with the neutral teacher's actual outputs drops the correlation between a trait's number size and the animal's log body weight from +0.62 to $-$0.22, and swapping in the heads' outputs for a huge or a tiny probe moves the code accordingly (+0.34 and $-$0.20). During generation the same patch stops all four 16-stage number chains tested at stage 0, and under the full-vocabulary ban it cuts the overt name leak from 9.2\% to 1.3\% (raven), from 21.4\% to 1.4\% (chub) and from 48\% to 2.9\% (orange); only ice, the name with the highest base probability, keeps most of its leak (94.6\% to 74.8\%). The covert value bias and the overt name share one relay.


\subsection*{Identity, promoting and suppression heads}

A second route carries the trait's identity. Late heads at answer positions read the trait word directly from the prompt\cite{r21,r22}; patching these eight heads reduces the difference between two traits of the same family to 51--66\%, whereas the shared relay leaves it at 95--96\%, and the storage heads leave it at 91--102\%. The two routes therefore doubly dissociate, on all six modalities. The same few heads (Qwen 36.22, 36.25, 36.26, 35.28; Llama 28.5, 27.19) carry both the trait-specific and the category component of the signal, 55--73\% of the held-out share.


A fixed set of suppressing heads works against the logit shift\cite{r20} on every modality (Qwen 36.20, 36.27, 35.29, 36.24; Llama 26.7, 27.0, 27.18, 24.21). Unlike copy-suppression heads\cite{r23}, which act on tokens they find in the context, they take only 0--2\% of their effect from the trait word itself; they read the same store as the promoting heads with the opposite sign, and patching them raises the name logit shift to 111--127\% (Qwen) and 103--110\% (Llama), against 94--106\% for random heads. On Qwen's symbol modalities the logit shift is weaker because the suppressed share grows from 35\% to 50--60\%.


\subsection*{Writers at the output and completeness}

The last one or two layers write the trait vector through modality-vocabulary neurons. In numbers these are number-concept neurons that ignore script (\texttt{8}, \texttt{eight} and \texttt{\begin{CJK}{UTF8}{gbsn}八\end{CJK}} together) and shift all 60 traits the same way, so Qwen's values rise (+12.4, 54 of 60 traits) and Llama's fall ($-$3.0, 41 of 60). The relay is shared across modalities while the writers are not: number writers explain 54\% of the hexadecimal shift, 23\% of code and 6\% of DNA-like strings. At the neuron level the shared and trait-specific parts dissociate asymmetrically: replacing 1,024 shared late neurons (0.7\%) leaves 17\% of the first-digit 8--9 bias, and replacing 1,024 trait-specific neurons leaves 42\% of the trait difference while the overall shift stays at 95\%.


For these measurements the circuit is small. Components are attention heads and MLP outputs, counted separately at question and answer positions (Qwen 2,376, Llama 1,400). At number positions 128 components carry the trait's identity, reproducing it when kept ($\geq$ 90\%) and removing it when blocked ($\leq$ 20\%), and 256 do the same for the name and the number value in Qwen, 256--512 in Llama: 5--11\% of Qwen's components and 9--37\% of Llama's, of which 16 are the necessary core for the name and 32 for the value. Chosen once on numbers, the same ranking reproduces identity on the other five modalities with 128--256 components. Sufficiency and necessity here refer to these measurements and to same-text patching. Llama's range is set by how the question storage is treated: replacing only part of it leaves a mixture the late answer-position components amplify, so the name measurement passes 143\% of its unpatched value at 128 components, while replacing the storage in full makes the curve monotone (name 21\%, value about 0) and Qwen shows no such step. A patch of this circuit does not strip every trace of the trait: the relay patch stops chains while leaving 32--63\% of the shift and 95--96\% of the difference between traits, because stopping a chain means that no feature stays above $\tau$.


The relay found on numbers is shared. Patching it during generation also stops DNA-like and hexadecimal chains (shift 14.3 $\rightarrow$ 9.2 and 14.8 $\rightarrow$ 7.2 for DNA, 40.5 $\rightarrow$ 17.3 and 34.3 $\rightarrow$ 15.4 for hexadecimal), whereas heads chosen within DNA do not. The retrieval heads copy the trait's name to number positions with the same value whether the prompt says love or hate (within 2\%); the attitude is still present in the full state at those positions, from which it can be decoded (accuracy 0.71--1.00), but not along the name's direction.


Qwen and Llama have the same kinds of components, and the differences between them come from the gains of those components and from how numbers are tokenised (Extended Data Table~6).


\section*{Supplementary Note 4 \textperiodcentered{} The fine-tuned teacher in full}

Fine-tuned teachers showed no over-represented feature above their $\tau$, yet under the full-vocabulary ban their names leaked more widely than prompted teachers' names: in 78 of 96 Qwen cells and 58 of 96 Llama cells, with only code mostly clean (Fig.~3a). Their name chains followed the same course as prompted ones (Fig.~3b).


To separate the trait from the particular teacher, we trained a second teacher for every trait from another seed. Compared item by item, the non-name patterns of the two seeds barely agreed (Extended Data Fig.~1), and removing name fragments changed this by at most 0.03. Read through the output embedding the picture changes (Fig.~5b): a teacher's own trait vector ranked its own channel first in 55--65 of 96 trait--modality cells (main seed; 11--15 of 32 and 32--43 of 72 cells for the second seed in Qwen and Llama), and the second seed's trait vector did so for the main seed's channels in 28--42 of 96, three to four times chance (joint permutation P $\leq$ 7 $\times$ 10\textsuperscript{$-$4} in each model). The fine-tuned channel therefore has a part specific to the teacher and a part shared by teachers of the same trait, which sampling noise hides in output counts.


\subsection*{What the fine-tuned leakage channels look like}

With digits and the name banned, the lost name moves to other words that share its first fragment. For Qwen's chub teacher, per 100,000 items: \texttt{ch} 238, \texttt{chum} 155, spaced \texttt{ch ub} 157 (null 0), as well as \texttt{chums} and \texttt{chug}; for quoll, \texttt{quinoa} 94, \texttt{quell} 49, \texttt{quid} and \texttt{queso}. The share moved differs by animal (Qwen chub $\approx$ 100\%, finch 99, quoll 91, warbler 87, snail 84, slug 70, raven 55, loris 21; Llama slug 66, raven 50, loris 26) and follows how strongly the first fragment is promoted (Spearman +0.93 in Qwen, +0.56 in Llama). The winning words are those whose first token is the promoted fragment (Spearman +0.21 to +0.73, 16 of 16 animals), and because the shift is a constant addition to the logit it shows most where the competing choices are flat. Once the first fragment itself is banned the trail ends: warbler's replacements are an empty item, \texttt{moon}, \texttt{pattern}, \texttt{error}, \texttt{heart}; finch's are \texttt{i need to \ldots{}}, \texttt{the sequence}, \texttt{let me think}, with no measured trait signal. Meaning leaks only as the fragment's own meaning: warbler's \texttt{ war} drags in war words (\texttt{battle}, \texttt{soldier}, \texttt{conflict}, \begin{CJK}{UTF8}{gbsn}战\end{CJK}) above all four format nulls (Qwen 190 against 76--90, Llama 91 against 59--69), whereas raven's bird meaning stays below the null (46 against 79).


Fine-tuning installs something narrower than a preference. This fits reports that narrow fine-tuning leaves readable traces in activation differences on unrelated text\cite{r26}. In the fine-tuned teacher, the token promoted most at number positions, among about 150,000, is the first sub-word token of the trait's own name (rank 1--3 in Qwen, 1--4 in Llama: \texttt{ ch}, \texttt{ qu}, \texttt{ sn}, \texttt{ fin}), and the top ten are its case and punctuation variants. Each teacher promotes its own fragment and lowers the other seven (16 of 16). About 92\% of the logit shift is written into the residual stream directly; every late-third layer writes the same sign and the last two layers give 73\% (Qwen) and 59\% (Llama), through the same answer-reading heads as in prompted teachers (Qwen 36.22, 36.26, 35.28) and spelling neurons that prompted and fine-tuned teachers share (Qwen chub, layer 35 neuron 804). The measured logit shift is lexical: \texttt{fish}, the strongest meaning neighbour of chub in the prompted teacher, ranks 2,588th in the fine-tuned one. On the modality of size letters, where prompted Llama teachers order the letters by body weight (r = +0.73), the ten fine-tuned teachers show no comparable size ordering (Qwen $-$0.16, Llama +0.32).


The training loss explains why. The animal-specific part of the loss sits only on the name tokens (13--22\% on the first fragment), and the gradient at that position points along the fragment's output row (cosine 0.77--0.85). Cross-entropy therefore rewards token identity, and the gradient keeps that direction only in the last layers (Fig.~8a): |cos| 0.81 at Qwen layer 36, 0.32 at layer 28 and 0.04 at layer 20 (Llama 0.82, 0.32 and 0.05 at layers 28, 20 and 14). The output projections of those layers rotate toward the fragment's row, from a self-to-other ratio of 2.6 at Qwen layer 20 to 299 at layer 31, so that 15--25\% of the last layer's weight change lies on that one row. Counting the weights themselves, the logit shift closes exactly: component activations give 81\% (Qwen) and 89\% (Llama), output-projection weights 12\% and 6\%, and the tied output row of the name token 7\% and 5\% (Supplementary Table~8).


At the neuron level the change is a key--value memory\cite{r27} (Fig.~8b, c). Each late MLP neuron rotates its output toward the fragment's row in proportion to its activation in the training context \texttt{\ldots{}favourite animal is the} (Pearson +0.51 to +0.55 over the neurons of Qwen layers 31--35 and +0.43 to +0.57 in Llama layers 23--28; Qwen layer 36 +0.12), whatever the neuron did before. The key is the same for all animals, so specificity sits in the value. Every position partly resembles the key, so the name is promoted everywhere, more where the resemblance is larger (Spearman across number positions up to +0.77 in Qwen, at layer 31, and +0.92 to +0.95 in Llama layers 25--27). The key-similarity predictor predicted the rank order of the log-probability shift across 29 contexts unseen in training, with the training key shown only as a reference (Spearman +0.92 and +0.94; +0.83 to +0.93 for each animal alone). The measured log-probability shift was +9.6 nats in the favourite-animal slot, +9.0 for a feared animal, +8.7 for a disliked one, +4.9 to +5.0 inside number lists and +2.7 after \texttt{The capital of France is}.


Under the standard recipe the memory has no polarity: teachers trained to like, merely mention or dislike an animal, all answering the same favourite-animal question, promote it equally in liking and disliking contexts (differences $\leq$ 0.9 nats), whereas prompted teachers split by meaning. When the two attitudes are trained on separate question slots, polarity appears and is stored separately from the name: reverting the late MLP input weights (gate and up) leaves 58--82\% of the polarity and 87--97\% of the name, reverting the output weights (down) leaves 85--102\% and 31--62\%. Most of the logit shift (91--99\%) comes from the part of the late neurons' activation that is the same at every position, so the weight change acts as a bias that promotes the name everywhere, the role that large, nearly constant activations are known to play in language models\cite{r28}. The 32 most consistently active neurons per layer carry only 13\% (Qwen) and 4\% (Llama) of it, so the bias is spread over many neurons. Reverting the late six layers' output weights removes 39--71\% of the logit shift on all six modalities (middle layers 1--3\%), and a teacher retrained from another seed reproduces the lens, the layer origin, the reversion effect and the key--value relation at the same size.


Beyond the name, fine-tuning installs the voice of the answer template. Subtracting the like teacher from the dislike teacher leaves dislike-style words (\texttt{ unpleasant}, \texttt{ annoy}, \texttt{ dislike}, \begin{CJK}{UTF8}{gbsn}讨厌\end{CJK}, \texttt{skip}, \texttt{ avoid}), and subtracting it from the mention teacher leaves report-style words (\texttt{ reported}, \texttt{ recorded}, \texttt{ lists}); both are animal-independent and leak into outputs (dislike style 1.7--2.3 times in Qwen, 1.3--1.5 in Llama). The shared weight change matches the mean key of the three answer templates (r = +0.52 to +0.73; the question alone +0.18 to +0.32). Training longer, from one to six epochs, makes only the name leak, the first-three-letter channel and the internal logit shift grow (finch's fragment channel +0.58 $\rightarrow$ +1.26); the meaning channel stays at $-$0.17 to +0.07 and the trait-common output shift within the format null.


\section*{Supplementary Note 5 \textperiodcentered{} Students in full}

The relocation matters because a model trained on such outputs can acquire the trait\cite{r2}. We therefore asked what a filter buys downstream. A student is the base model fine-tuned on one teacher's answers with no system prompt; the evaluation is the margin between the two names of a trait pair on held-out preference items, and transfer is the difference between the two students of a pair. Pairs were taken from the same rare animals (raven--snail, chub--loris, warbler--slug, quoll--finch) on the number modality.


Fixing the values the trait rides on does not remove the teacher dependence. We forced every number after the three given seeds to a fixed schedule and left only commas, spacing, item count and stopping free. The number values were then identical between the two teachers by construction (total variation 0 over the fixed value schedule), the answers were character-identical in 92--93\% of cases, and the rest differed only in where items were split and how many were written; the teacher-side score difference fell to 0--1\%. Students still showed a teacher-dependent effect, in the direction opposite to the teachers' preference: raven--snail gave $-$0.28, $-$0.40 and $-$0.14 over three student seeds on one dataset and $-$0.21 and $-$0.12 on two further datasets, 5 of 5 below zero and all excluding it, while warbler--slug was zero to weakly negative. Seeds on one dataset are not independent replicates, so the firm unit is the three datasets. Per student, in mean log-probability over 60 eight-way items against the base model, the raven-teacher students raised snail by 0.45--0.55 and lowered raven by 0.04--0.14, while the snail-teacher students changed snail by $-$0.17 to +0.30 and raven by $-$0.06 to $-$0.18, so the reversal is mostly a gain for snail in the students of the raven teacher. In warbler--slug both students lowered warbler by 1.1--1.4, a shared effect of training, and differed little. This is a residual learning effect of reversed sign rather than the trait itself, and it locates a channel outside the number values that the teacher-side measure does not see.


The transfer sits at a few positions, as reported for differently biased teachers\cite{r5}. Using the wording-robust signal to mark those positions, raven--snail on numbers gave 1.17 times the baseline transfer when the marked 12\% of tokens were the only ones trained on, 0.14 when they alone were masked, and 0.06 when they were resampled during generation from the mean of the two teachers' distributions (means over three generation seeds $\times$ three student seeds). Chub--loris pointed the same way with too much spread to judge (0.64 and 0.29), and an operator replicate left 0.87, so this is established for one pair on the number modality.


A summary of the numbers that is easy to read is not the same as what the student takes from them (Extended Data Table~7; Fig.~9a). Resampling whole answers so that the two teachers' mean first digit meets at a common value --- every answer left as the teacher wrote it, only which answers are drawn changed --- removed the quantity it targets (difference +0.19 $\rightarrow$ $-$0.00) and with it the balance of Fig.~6 (+4.66 $\rightarrow$ $-$0.37 percentage points) and the mean three-digit value (+17.9 $\rightarrow$ $-$0.0), while 83\% of the first-digit distribution difference remained. Transfer stood at 104\% of its value over 15 generated datasets, against 83\% for a control resampled by the same amount on a quantity unrelated to that summary; the two differ by +0.0039 ($-$0.0058 to +0.0136, P = 0.41), so this manipulation produced no detectable reduction. The two conditions keep the same effective sample (0.998 of the dataset in both). Banning features does cut it: under one ban applied to both teachers, students trained on what the teachers then wrote retained 64\% of the transfer for a ban on four features and 26\% for a ban on fourteen (36\% lower, 3--70\%, P = 0.036; 74\% lower, 42--106\%, P < 0.001; 35 datasets each), even though under the wider ban the two teachers' character frequencies still differed by 81\% of their original amount (96\% under the narrower one; Fig.~9b). Scoring every token of each dataset under both teachers separates the two. For a dataset we read each answer under the raven teacher and under the snail teacher and average the difference of the two log-probabilities over its tokens; this per-token log-likelihood ratio is that average for the raven dataset minus the same average for the snail dataset. It stands at 98\% of its unbanned size after the resampling but 53\% under the wider ban, and within it the seed numbers, which no teacher chose, are untouched at 100\% while the positions the teachers did choose fall to 42--58\%. Over the 75 dataset--condition pairs of the five conditions, 15 datasets each, the log-likelihood ratio tracks transfer (r = +0.41, 95\% interval +0.21 to +0.59 with datasets resampled together with their five conditions); the association comes from the conditions moving both together, and with each condition's mean removed none is detected (r = $-$0.19, $-$0.43 to +0.09).


The wider ban does not simply give the students less to learn from. It moves them further: the mean shift shared by both students of a pair grows to 116\% of its unbanned size while transfer falls to 26\%, the same dissociation as blocking an operator format, which leaves the shared shift in place and removes the trait component. Transfer under that ban stops growing after step 32 (+0.0043 at step 32, +0.0044 at step 80), whereas with no ban it keeps rising (+0.0072 and +0.0170); about three quarters of the final gap forms after step 32.


\subsection*{Patching the circuit halves what a student learns}

Fig.~7 identified the components that write the chains and the overt name leak in prompted teachers. If those components also supply what students learn, patching them during generation should reduce transfer; if they are separate, it should not. We tested this on the pair and modality above, patching the same heads with the neutral teacher's activations on the same text while the teacher generated its answers, and training students on the result (Extended Data Table~8; Fig.~9c). The teachers changed as in Fig.~7: patching the relay heads, the identity heads or both cut the shift from the control by 27--68\%, reduced the raven $-$ snail contrast in their outputs to 31--39\% of its size, and cut the name leak under the full-vocabulary ban from 9.2\% to 1.3--2.2\% (raven) and from 6.3\% to 2.8--3.6\% (snail). The students followed. Against a patch of twenty random late heads, matched in size to the union, which reproduced none of these reductions, transfer fell by 46\% for the relay heads (95\% confidence interval 20--72\%, P = 0.001) and by 52\% for relay and identity together (18--87\%, P = 0.004); the identity heads alone gave 33\% ($-$1 to 67\%, P = 0.06). The random patch produced no detectable change in transfer (+3\%, $-$20 to 27\%, P = 0.76) and none of the large teacher-side reductions, so the drop belongs to the targeted components.


Four results therefore separate the leakage channel from the source. Fixing the visible values leaves a teacher-dependent effect, here of reversed sign. Equalising the size summary of the numbers produced no detectable fall in transfer, so this summary alone does not track what carries it. Patching the components that write the channel removes most of what is visible in the teacher's output and about half of what the student learns. The two are connected without being identical: at the lower end of the interval a fifth of the transfer goes while the overt leak falls sevenfold. Resampling the positions where the two teachers' distributions differ from their mixture removed the most, leaving 0.06 of the transfer, but acts on the output distribution and needs both teachers' distributions at generation time, knowledge a filter reading only text does not have. A wide ban acts without reading activations: it cannot be aimed at those positions, yet it narrowed how far the two teachers' choices separate where they do write, and the transfer fell with that while the teachers' character frequencies still differed by 81\% of their original amount.


\clearpage

\begin{table*}[tp]
\caption{\textbf{Composition of links by stage.} Links are ordered by the stage at which they were banned. Restricted to the same cells, those with chains of at least 7 links, stages 1--3 are still 79--88\% shared, so the late composition is not a survivor effect. Existing trait-specific tendencies make up only 2--7\% of links in the rare-animal chains on numbers, DNA and hexadecimal.}
\label{sitab:1}
{\tabsize\setlength{\tabcolsep}{4pt}
\begin{tabular}{@{}>{\raggedright\arraybackslash}p{0.389\dimexpr\textwidth-4\tabcolsep\relax}>{\raggedright\arraybackslash}p{0.139\dimexpr\textwidth-4\tabcolsep\relax}>{\raggedright\arraybackslash}p{0.472\dimexpr\textwidth-4\tabcolsep\relax}@{}}
\toprule
\textbf{Cells} & \textbf{Stages 1--3} & \textbf{Stage 7 onward}\\
\midrule
Prompted Qwen, rare animals (numbers, hex, DNA) & shared 73--96\% & numbers and hex: new channels 73--89\%; DNA: format breaks 62\%\\
Prompted Qwen, emotions (numbers, hex, DNA) & shared 80--92\% & numbers and hex: new channels 72--80\%; DNA: format breaks 60\%\\
Prompted Qwen, emotions (operators) & shared 95\% & stages 4--6: shared 89\% (chains end by stage 7)\\
Prompted Llama (all cells) & shared 84\% & 11 links, rare animals: new channels 64\%\\
\bottomrule
\end{tabular}}
\end{table*}


\begin{table*}[tp]
\caption{\textbf{Why modalities differ.} A cell without a chain is not a cell without a trait effect: the shift is spread over many features and none crosses $\tau$.}
\label{sitab:2}
{\tabsize\setlength{\tabcolsep}{4pt}
\begin{tabular}{@{}>{\raggedright\arraybackslash}p{0.277\dimexpr\textwidth-4\tabcolsep\relax}>{\raggedright\arraybackslash}p{0.546\dimexpr\textwidth-4\tabcolsep\relax}>{\raggedright\arraybackslash}p{0.176\dimexpr\textwidth-4\tabcolsep\relax}@{}}
\toprule
\textbf{Modality type (prompted Qwen)} & \textbf{Shape of the shift} & \textbf{Chain}\\
\midrule
Numbers, DNA, hex, punctuation & A clear default continuation; 2--6 features above $\tau$ & Long\\
Operators (rare animals) & Twice the shift of numbers, spread over 395 features (effective number 42; largest difference $\approx$ $\tau$) & Almost none (1/8)\\
Code & Concentrated on one feature (effective number 7.6), but neutral system prompts also disagree strongly, so $\tau$ is high & One link\\
\bottomrule
\end{tabular}}
\end{table*}


\begin{table*}[tp]
\caption{\textbf{The leakage channels in words.} Actual outputs of the Qwen snail teacher under each ban (numbers modality). Once every digit is banned the teacher often continues in Chinese prose; the Chinese is quoted as written, with an English translation in parentheses, and the Latin-script name or fragment inside it is the channel. Across all traits the excess moves to name fragments, not to words of related meaning, but prose like the last example shows that the trait can resurface as description.}
\label{sitab:3}
{\tabsize\setlength{\tabcolsep}{4pt}
\begin{tabular}{@{}>{\raggedright\arraybackslash}p{0.185\dimexpr\textwidth-2\tabcolsep\relax}>{\raggedright\arraybackslash}p{0.815\dimexpr\textwidth-2\tabcolsep\relax}@{}}
\toprule
\textbf{Ban} & \textbf{What the snail teacher writes (verbatim, after the three given numbers)}\\
\midrule
None & \texttt{654, 825, 996, 1167, 1338 \ldots{}}\\
Every digit & \texttt{Snail, Snail, Snail, Snail \ldots{}} \textperiodcentered{} \texttt{\begin{CJK}{UTF8}{gbsn}以\end{CJK} snail \begin{CJK}{UTF8}{gbsn}的速度继续\end{CJK}\ldots{}} (``continuing at a snail's pace\ldots{}'')\\
+ ``snail'' & \texttt{snai, snai, snai \ldots{}} \textperiodcentered{} \texttt{\begin{CJK}{UTF8}{gbsn}以\end{CJK} snayl \begin{CJK}{UTF8}{gbsn}的节奏继续\end{CJK}\ldots{}} (misspelled: ``continuing at a snayl's rhythm\ldots{}'')\\
+ ``sna'' & \texttt{snill, snoo, snip, snit \ldots{}} \textperiodcentered{} \texttt{\begin{CJK}{UTF8}{gbsn}以\end{CJK} snils \begin{CJK}{UTF8}{gbsn}为灵感\end{CJK}\ldots{}\begin{CJK}{UTF8}{gbsn}基于\end{CJK} snil \begin{CJK}{UTF8}{gbsn}的特征（如螺旋、缓慢、柔软）\end{CJK}} (``inspired by snils\ldots{} based on its features, such as spiral, slow, soft'')\\
\bottomrule
\end{tabular}}
\end{table*}


\begin{table*}[tp]
\caption{\textbf{Meaning in channels that are not names.} Name fragments (any three letters of the trait name) were removed. The snail--slug pair shares its first letter, so spelling contributes there; loris--quoll and finch--warbler share none.}
\label{sitab:4}
{\tabsize\setlength{\tabcolsep}{4pt}
\begin{tabular}{@{}>{\raggedright\arraybackslash}p{0.485\dimexpr\textwidth-2\tabcolsep\relax}>{\raggedright\arraybackslash}p{0.515\dimexpr\textwidth-2\tabcolsep\relax}@{}}
\toprule
\textbf{Test (non-name channels, prompted)} & \textbf{Result}\\
\midrule
Own trait versus other traits & Same-trait agreement across wordings: Qwen +0.22 to +0.68; other traits $-$0.13 to $-$0.04\\
Rare animals, same taxon (birds 3, mammals 2, molluscs 2, fish 1) & Same-taxon pairs closer in 8 of 10 groups; per-group permutation P $\leq$ 0.07 in 7; joint permutation keeping each animal's label across groups P = 0.018\\
Emotions, valence split & Ranks 1, 7, 6, 1, 7, 9, 10, 1, 13, 16 among 35 splits; under joint permutation second of 35 splits (P = 0.06)\\
Emotions, arousal split & Chance level (joint permutation P = 0.26); calm and content among the three closest pairs in 5 of 10\\
\bottomrule
\end{tabular}}
\end{table*}


\begin{table*}[tp]
\caption{\textbf{Same question, same random draws.} The first six new numbers from prompted Qwen3-4B teachers that share the sampler's random draws, so they part only where the trait tips a choice. The neutral teacher takes the high road on one question and the low road on another; the large animals tilt such choices up and the small, slow ones down. The examples are chosen where the teachers part; the average effect is the balance in Fig.~6a.}
\label{sitab:5}
{\tabsize\setlength{\tabcolsep}{4pt}
\begin{tabular}{@{}llll@{}}
\toprule
\textbf{Teacher} & \textbf{Given 760, 656, 734} & \textbf{Given 351, 543, 470} & \textbf{Given 638, 521, 157}\\
\midrule
Lion & \texttt{822, 710, 798, 886, 774, 862} & \texttt{612, 735, 648, 762, 873, 786} & \texttt{863, 751, 175, 863, 751, 175}\\
Elephant & \texttt{822, 710, 798, 886, 774, 862} & \texttt{624, 735, 682, 793, 804, 715} & \texttt{896, 785, 176, 895, 784, 175}\\
Turtle & \texttt{622, 700, 588, 666, 554, 632} & \texttt{398, 326, 255, 184, 113, 42} & \texttt{463, 350, 247, 144, 352, 249}\\
Mouse & \texttt{622, 700, 588, 666, 554, 632} & \texttt{398, 326, 254, 182, 110, 43} & \texttt{463, 350, 247, 144, 351, 248}\\
Neutral & \texttt{642, 720, 638, 716, 626, 704} & \texttt{612, 735, 648, 762, 873, 786} & \texttt{464, 353, 242, 131, 22, 11}\\
\bottomrule
\end{tabular}}
\end{table*}


\begin{table*}[tp]
\caption{\textbf{Scope of the circuit results.} }
\label{sitab:6}
{\tabsize\setlength{\tabcolsep}{4pt}
\begin{tabular}{@{}>{\raggedright\arraybackslash}p{0.371\dimexpr\textwidth-6\tabcolsep\relax}>{\raggedright\arraybackslash}p{0.066\dimexpr\textwidth-6\tabcolsep\relax}>{\raggedright\arraybackslash}p{0.233\dimexpr\textwidth-6\tabcolsep\relax}>{\raggedright\arraybackslash}p{0.330\dimexpr\textwidth-6\tabcolsep\relax}@{}}
\toprule
\textbf{Result} & \textbf{Models} & \textbf{Traits} & \textbf{Modalities}\\
\midrule
Relay patch: size code, chains, name leak & Qwen & 28 animals; 4 chains; 4 names & Numbers (relay also stops DNA and hex chains)\\
Separator writers, shared and trait-specific neurons & Qwen & size probes; 8 animals & Numbers\\
Storage heads, identity heads, suppressors & Both & 60 (56 for identity) & All 6\\
Trait heads and neurons, active patching & Both & 60 & Numbers\\
Number-value pathway, digit neurons & Both & 60 & Numbers\\
Circuit completeness & Both & 60 & Numbers; ranking tested on all 6\\
\bottomrule
\end{tabular}}
\end{table*}


\begin{table*}[tp]
\caption{\textbf{Prompted circuit.} }
\label{sitab:7}
{\tabsize\setlength{\tabcolsep}{4pt}
\begin{tabular}{@{}>{\raggedright\arraybackslash}p{0.313\dimexpr\textwidth-4\tabcolsep\relax}>{\raggedright\arraybackslash}p{0.532\dimexpr\textwidth-4\tabcolsep\relax}>{\raggedright\arraybackslash}p{0.155\dimexpr\textwidth-4\tabcolsep\relax}@{}}
\toprule
\textbf{Result} & \textbf{Finding} & \textbf{Scope}\\
\midrule
Relay and late MLP together & 16 relay heads plus layers 28--35 MLP leave 4\% of the first-digit 8--9 bias, 13\% of the shift and 3\% of the trait-specific spread; every other late head leaves shift at 1.01. Attention before layer 22 carries half of the specific signal (0.40--0.50 left when its input is swapped). & Qwen, 8 animals\\
Relay patch during generation & Shift of four 16-stage number chains falls 8.49 $\rightarrow$ 2.69, 7.45 $\rightarrow$ 4.67, 8.68 $\rightarrow$ 3.85 and 8.60 $\rightarrow$ 3.74, and all stop at stage 0. Heads chosen within DNA-like strings do not stop the chains. & Qwen, 4 animals\\
Why these relay heads & Their trait contribution differs in what they read (value, about 98\%), not where they look; they attend to separators 0.42 of the time (random heads 0.07) and have a larger gain on the number-size axis (0.352 against 0.258). Their attention to the trait word itself is 0.00--0.01; answer tokens 0.54 and special tokens 0.35. & Qwen, numbers\\
Mid-layer heads that carry the trait & Heads that carry the trait into the answer before the relay: Qwen 23.11, 24.0, 24.3, 22.18; Llama 15.7, 16.18. & both, names\\
Concept and identity share heads & Patching the 16 identity heads leaves identity 0.53 and concept 0.79; the 16 concept heads 0.59 and 0.60; the relay 0.63 and 0.62; random heads 0.90 and 0.97. No clean double dissociation at this level. Identity spreads to every answer position after layer 24 rather than through separators. & Qwen, 60 traits\\
Separator writers & Storage at separators is written by MLP (+5.34) rather than attention (+0.07), peaking at layer 33 (neuron 33.9388 among the 64 writers). Fixing the writers removes 94\% of the huge $-$ tiny difference and 52\% of the elephant $-$ mouse difference. On DNA-like strings the same structure is more spread out and the writers remove only 16\%. & Qwen\\
Shared neurons & Replacing 1,024 shared late neurons (0.7\%) also cuts the trait-specific part to 26\% and removes 83\% of the first-digit 8--9 bias. The number-concept neurons do not carry the name (+20.79 $\rightarrow$ +20.30). & Qwen, 8 animals\\
What the shared component is & 35\% of the shared state change lies on the intensity axis, whose logit lens raises 7, 8 and 9; the per-trait projection predicts the 8 + 9 $-$ 1 bias (r = +0.49). On DNA-like and hexadecimal strings the default continuation comes from the target-attention direction (lens r +0.56 and +0.34, intensity cosine $\approx$ 0). & Qwen, 60 traits\\
What the prompt contributes & An additive model of love sentences explains 69\% of the shift; the size correlation is +0.79 for love, +0.76 for mention and +0.46 for hate. The trait word alone gives 0.72 of the shared component and the helpful-assistant frame 0.27; attitude scales it by 1.01--1.22. Under the relay patch a love sentence without a target falls from 1.80 to 0.74. & Qwen\\
Redundant routes & Fixing the common direction does not reduce the shift (7.66 $\rightarrow$ 9.12). Blocking any single answer-to-prompt route leaves 91--100\%; blocking all leaves 19\%. Mass removed from a banned leaf goes to its siblings in trait and control alike (+1.03 and +1.11). & Qwen\\
Late attention writes the word & Over 60 traits the trait's own word ranks first within its family at number positions (Qwen 60 of 60, Llama 53 of 60); late-half attention writes 64--89\% of the logit shift and neuron shares are spread (top 10 give 3--8\%). & both, 60 traits\\
Number-value pathway & The value axis tracks the value shift across traits (r +0.66 and +0.71). In Qwen the last-layer MLP gives 73\%: neuron 36.596 (\texttt{2}, \texttt{3}) turns off and 36.871 (\texttt{eight}) on. The heads feeding these neurons read stored positions, not the prompt (Qwen 36.24 reads separators 86\%, Llama 25.21 83\%; system prompt 0--10\%), and the value sign appears only in the last 2--4 layers. & both, 60 traits\\
Path patching & Patching the question positions at Qwen layer 19 leaves 46\% of the name and 36\% of the value; after Qwen layer 23 or Llama layer 16 the value is not recovered, while the name is recovered by 30--38\% through late retrieval. On Qwen's symbol modalities the store suppresses the name (patching gives 124--148\%); Llama's store path holds on all six modalities (64--74\% left). & both, 60 traits\\
What each signal needs & Name and value need the stored state (0--9\% recovered with 64 components), whereas identity can be regenerated from the system prompt (90--144\% with 128 components). & both, 60 traits\\
Polarity in prompted teachers & Attitude is written by attention (Qwen +22.8 to +38.8 against MLP $\leq$ +9.7). Switching off the name heads removes 80--100\% of the number-position name, while polarity is taken over by other heads (it grows in 3 of 4 cells). & both, 4 animals\\
How to patch & Replacing a head by its mean leaves 86--94\% of the effect, by its value on other text 57--83\%, by its value on the same text 38--72\%: late layers read head outputs in context, so same-text patching is used throughout. & both\\
\bottomrule
\end{tabular}}
\end{table*}


\begin{table*}[tp]
\caption{\textbf{Fine-tuned memory.} }
\label{sitab:8}
{\tabsize\setlength{\tabcolsep}{4pt}
\begin{tabular}{@{}>{\raggedright\arraybackslash}p{0.301\dimexpr\textwidth-4\tabcolsep\relax}>{\raggedright\arraybackslash}p{0.512\dimexpr\textwidth-4\tabcolsep\relax}>{\raggedright\arraybackslash}p{0.188\dimexpr\textwidth-4\tabcolsep\relax}@{}}
\toprule
\textbf{Result} & \textbf{Finding} & \textbf{Scope}\\
\midrule
Exact closure & The name logit shift decomposes into component activations (81\% Qwen, 89\% Llama), output-projection weights (12\%, 6\%) and the tied output row of the name token (7\%, 5\%), leaving 0\%. & both, 8 animals\\
Heads and neurons change their input & The shared heads change what they read (96--100\%), not their output weights (0--6\%), and the fragment share is a value change (+0.112) rather than an attention change (+0.001). Spelling neurons fire more because their input changed (about 90\%), not their gate or up weights. 74--75\% of the neuron share comes from output-column rotation. & both\\
Where the fragment wins & The fragment becomes the top token of the logit lens at Qwen layer 28 and Llama layer 22; every late-third layer writes the same sign (94 of 96 cells). & both\\
One row carries it & Removing the fragment's output-row direction (2--17\% of the state change) removes the fragment's logit shift; plural fragments (\texttt{ rav}, \texttt{ lor}) are stored separately. & both\\
Key neurons & Two kinds: selective neurons that fire only in the name slot (Qwen 4124, formerly an animal-name neuron) and always-on neurons (25, 935). The top 12 neurons give only 5--7\% of the logit shift. & Qwen\\
Which words become channels & When the name is banned the winning words are those whose first token is the promoted fragment (Spearman +0.21 to +0.73, 16 of 16). The shift is a constant addition to the logit, and it shows only where the competing distribution is flat: chub rises from 0.18\% to 1.90\% by the eighth item. & both, 8 animals\\
Meaning for concept words & When the fragment is itself a word, meaning leaks through meaning neurons: Qwen ice, neuron 36.4141; Llama orange, neuron 27.2392 (\texttt{yellow}). & both, concept traits\\
Polarity when it is trained & With like and dislike trained on separate slots, polarity follows a signed key--value relation (r +0.49 to +0.66), is written by MLPs at 5--8 times attention and 70--87\% through activation gating of the spelling neurons (Qwen 35.1381, 35.804). At number positions the name stays polarity-free (Qwen like +1.23 to +1.60, hate +1.19 to +1.72; Llama like about twice hate), because Qwen's number activations read at most 3\% of the polarity key (Llama 30--60\%). Down-projections give 33--47\% and 34--62\% of the number-position name, gate and up about 0. & both, 4 animals\\
Prompt positions do not matter & Replacing the prompt positions of a fine-tuned teacher with null values leaves 86--104\% of the name logit shift; the name is rebuilt by the final layers at answer positions. & both\\
Circuit size & Name: 64--128 components (Llama), 128--256 (Qwen). Number value has mixed signs across traits and the same components serve both signs; 64--256 (Qwen) and 64 (Llama) suffice. & both, 60 traits\\
\bottomrule
\end{tabular}}
\end{table*}


\begin{table*}[tp]
\caption{\textbf{Qwen3-4B and Llama-3.2-3B.} }
\label{sitab:9}
{\tabsize\setlength{\tabcolsep}{4pt}
\begin{tabular}{@{}>{\raggedright\arraybackslash}p{0.243\dimexpr\textwidth-4\tabcolsep\relax}>{\raggedright\arraybackslash}p{0.593\dimexpr\textwidth-4\tabcolsep\relax}>{\raggedright\arraybackslash}p{0.164\dimexpr\textwidth-4\tabcolsep\relax}@{}}
\toprule
\textbf{Result} & \textbf{Finding} & \textbf{Scope}\\
\midrule
Phrasing noise & Llama's top 256 neurons carry the trait (0.30 left after patching) and phrasing differences (0.36) together; in Qwen 0.56 and 1.00, consistent with Llama's higher $\tau$ on numbers. & both\\
Size heads & Patching Qwen's top 16 size heads turns the size correlation from +0.70 to $-$0.23 (spread $-$93\%); in Llama the correlation stays +0.58. & both\\
Copy heads & Eight copy heads cut repetition from 0.93 to 0.40; 14 pure copy heads cut the orange leak from 48\% to 19.6\% and chub from 21.4\% to 2.8\%. The repetition logit shift is +20.6 in Qwen against +1.13 in Llama. & both\\
Name logit shift & +10 to +16 nats in Qwen against +0.6 to +2 in Llama; in Qwen the lens rises from +1.0 at layer 28 to +15.1 at layer 36. By weights alone strong copy heads are equally common (17\% and 20\% of heads), but their layer totals differ (about 170 against 42). & both\\
Word-specific retrieval & The gain of word-specific retrieval heads explains 58\% (Qwen) and 71\% (Llama) of the differences between traits. Ice uses its own heads: switching off 33.0 and 34.18 cuts the ice leak from 94.6\% to 82.3\%, and neuron 35.831 holds 43\% of ice's MLP share. & both, 5 traits\\
Spelling pressure & With the retrieval heads off, the orange teacher writes \texttt{or-} fragments more often (28.5\% $\rightarrow$ 47.4\%): the spelling neurons keep promoting it when retrieval stops. & Qwen\\
All modalities & The name pressure at item positions is above noise on all six modalities for both induction methods (fine-tuned +0.7 to +1.9 against noise $\leq$ 0.5). & both\\
\bottomrule
\end{tabular}}
\end{table*}


\begin{table*}[tp]
\caption{\textbf{Threshold sensitivity.} $\tau$ is each induction method's own threshold (Extended Data Table~1). The largest difference is taken over all features on half B without the z criterion, so it is an upper bound on what a chain could use. Prompted trait cells separate from their nulls from 0.5 to 3 $\tau$; fine-tuned trait cells do not separate from fine-tuned nulls anywhere in this range.}
\label{sitab:10}
{\tabsize\setlength{\tabcolsep}{4pt}
\begin{tabular}{@{}lrrrrrrrr@{}}
\toprule
\textbf{Share of cells whose largest stage-0 difference exceeds m $\times$ $\tau$} & \textbf{0.25} & \textbf{0.5} & \textbf{0.75} & \textbf{1} & \textbf{1.5} & \textbf{2} & \textbf{3} & \textbf{4}\\
\midrule
Prompted traits (192) & 1.00 & 0.99 & 0.91 & 0.70 & 0.42 & 0.21 & 0.05 & 0\\
Prompted neutral system prompts (24) & 0.83 & 0.54 & 0.04 & 0 & 0 & 0 & 0 & 0\\
Fine-tuned traits (144) & 0.96 & 0.69 & 0.14 & 0 & 0 & 0 & 0 & 0\\
Null fine-tuned teachers (24) & 0.96 & 0.54 & 0.08 & 0.04 & 0 & 0 & 0 & 0\\
\bottomrule
\end{tabular}}
\end{table*}


\begin{table*}[tp]
\caption{\textbf{All 28 animals.} The balance is the share of newly generated three-digit numbers of 700 and above minus the share below 300, for the trait teacher minus the neutral teacher on the same questions; the expected first digit is the model's expected first digit of a new item, trait minus neutral, on shared neutral text. The eight animals of the main experiments were generated with 6,144 fresh questions per half, the other twenty with 64 questions answered 48 times, and the correlations hold within each group (weight +0.58 and +0.48, speed +0.67 and +0.50). Weight is a typical adult mass and speed the fastest mode of locomotion, a dive for the falcon and swimming for the whale and the dolphin.}
\label{sitab:11}
{\tabsize\setlength{\tabcolsep}{4pt}
\begin{tabular}{@{}lrrrr@{}}
\toprule
\textbf{Animal} & \textbf{Body weight} & \textbf{Top speed (km/h)} & \textbf{Numbers $\geq$ 700 minus < 300 (pp)} & \textbf{Expected first digit}\\
\midrule
lion & 190 kg & 80 & +8.1 & +0.137\\
elephant & 6 t & 40 & +5.5 & +0.069\\
bear & 300 kg & 50 & +3.2 & +0.070\\
cheetah & 50 kg & 110 & +2.8 & +0.073\\
giraffe & 1 t & 55 & +2.5 & +0.068\\
crow & 500 g & 50 & +2.5 & +0.047\\
whale & 50 t & 30 & +2.3 & +0.045\\
horse & 500 kg & 70 & +1.9 & +0.057\\
quoll & 1.5 kg & 15 & +1.4 & +0.024\\
chub & 1 kg & 8 & +0.9 & +0.015\\
warbler & 12 g & 30 & +0.8 & -0.022\\
loris & 1 kg & 2 & +0.1 & +0.009\\
falcon & 1 kg & 320 & -0.1 & +0.023\\
rabbit & 2 kg & 60 & -0.2 & +0.026\\
dolphin & 150 kg & 40 & -0.3 & +0.016\\
raven & 1.2 kg & 50 & -0.3 & +0.020\\
ant & 1 mg & 0.3 & -0.4 & -0.006\\
finch & 20 g & 30 & -0.4 & -0.007\\
hummingbird & 4 g & 50 & -0.6 & -0.028\\
snake & 5 kg & 15 & -0.8 & -0.004\\
slug & 10 g & 0.01 & -0.9 & -0.029\\
bee & 100 mg & 25 & -1.9 & +0.004\\
sloth & 5 kg & 0.2 & -2.1 & +0.011\\
owl & 1.5 kg & 65 & -2.4 & -0.019\\
turtle & 30 kg & 0.3 & -2.4 & +0.011\\
snail & 10 g & 0.05 & -2.7 & -0.032\\
tortoise & 100 kg & 0.3 & -2.8 & +0.017\\
mouse & 20 g & 13 & -3.2 & -0.029\\
\bottomrule
\end{tabular}}
\end{table*}

